\section{静电场中的导体}

\subsection{导体的基本性质}
导体是一类导电能力极强的物质，为了研究导体和静电场的相互作用，我们就必须依据导体的属性为其建立一个基本模型，为了讨论方便，下面我们以金属导体为例。

在金属导体中，由于原子中最外层的电子与原子核之间的吸引力较弱，因此常常容易摆脱原子的束缚在各个原子之间自由运动，这种电子称为\uwave{自由电子}。尽管自由电子可以在金属内不受阻力的运动，但是自由电子若要脱离金属则会受到极大的阻力，我们将自由电子脱离金属克服阻力所需做的功称为金属的\uwave{逸出功}，不同金属的逸出功不同，通常用\uwave{电子伏}作为其单位，记作~\si{eV}，当一个电子经过电势差为$1$\si{V}的电场空间时电场力所做的功就为$1$\si{eV}。
\begin{equation*}
    1\si{eV}=1.60\times 10^{-19}\si{C}\times 1\si{V}=1.60\times 10^{-19}\si{J}
\end{equation*}
由此可见~\si{eV}~即元电荷$e=1.60\times 10^{-19}\si{C}$乘以伏特~\si{V}，注意电子伏~\si{eV}~是一个能量单位。

\subsection{导体的静电平衡条件}
在金属导体中，金属中的正离子按照一定的方式有规则的排列为晶格，金属导体不带电时，自由电子的负电荷和晶格的正电荷数值相等，两者的电效应相互抵消，从而呈电中性，此时导体中的自由电子只作微观的热运动，而没有任何规则运动，因此也形象的称为\uwave{电子气}。

现在我们来研究，不带电的导体置于静电场中的情况。

如图\Refx{图}{导体的静电平衡}所示，当不带电的导体置于外电场$\vb*{E}_0$时，由于导体中包含大量可以自由运动的电子，从而在外电场的作用下，电子在电场力方向上定向运动，但是由于电子从外电场获得的能量通常不足以克服逃逸所需的金属的逸出功，电子将大量聚集在导体的表面附近，从而，导体的一侧带负电荷，导体的另一侧可以认为相应的带正电荷。
\begin{Figure}{导体的静电平衡}
    \subfigure[静电平衡前]
    {
        \inputfigure{Chapter09A_Figure01}
    }\qquad\qquad
    \subfigure[静电平衡后]
    {
        \inputfigure{Chapter09A_Figure02}
    }
\end{Figure}
我们将这种在外电场的作用下，使得原先并不带电的导体两侧分别带上等量异号的电荷的现象，称为\uwave{静电感应}。由于静电感应的作用，积聚在导体两侧的正负电荷间将产生一个附加的反向电场$\vb*{E}$，此时导体内的总场强为$\vb*{E}+\vb*{E}_0$，随着电子定向漂移运动的持续进行，反向电场的场强逐渐增大至外电场，此时导体内的总场强为$\vb*{E}+\vb*{E}_0=\vb*{0}$，此时电子的定向漂移运动停止，导体上的电荷分布达到稳定状态，这种稳定状态称为\uwave{导体的静电平衡}。

事实上，对于一般形状的导体，无论其是否带电，当达到静电平衡时
\begin{enumerate}
    \item \textbf{导体内的电场强度处处为零}。
    \item \textbf{导体表面的电场强度与表面垂直}。
\end{enumerate}
对于第一点，如果达到静电平衡时导体内某处的场强不为零，那么该处的电子将在电场力的作用下做定向运动，这就破坏了静电平衡的稳态，这一点在上面的例子中已经充分反映了。

对于第二点，如果达到静电平衡时导体表面某处的场强并不与表面垂直，那么该处的电子将在电场力的切向分量作用下沿导体表面发生漂移，从而破坏了静电平衡的稳态。当表面电场与表面垂直时电荷分布是稳定的，因为尽管此时表面电子受到垂直向外的电场力，但金属的逸出功同时对这些电子形成了一个反向的阻力，因此其所受的合力仍为零。

由以上两点可以得到一个重要推论，当导体达到静电平衡时，\textbf{导体是一个等势体，即导体内和导体表面的电势处处相等}，这是因为导体内部的场强处处为零，即$\vb*{E}=\vb*{0}$
\begin{equation*}
    U_{ab}=\Int[A][B]{\vb*{E}\cdot\dif{\vb*{r}}}=0
\end{equation*}
因此任意两点间的电势差也为零，故导体的电势处处相等。

导体达到静电平衡的过程是极为复杂的，但是实验表明，导体达到静电平衡的时间是极为短暂的，这一时间大约在$1\times 10^{-6}$\si{s}的数量级，因此在接下来的内容中，并不会关注“导体如何达到静电平衡”，而是会考虑“导体达到静电平衡后的性质”这样的问题。

\subsection{导体在静电平衡时的电荷分布}
在上一小节中我们讨论了静电平衡时导体的电场强度和电势，本小节将研究静电平衡时导体的电荷分布，这可以从\textbf{导体内部的电荷分布}和\textbf{导体表面的电荷分布}两个方面来考察。
\begin{Figure}{静电平衡时的电荷分布}
    \subfigure[导体内部的电荷分布]
    {
        \label{图:导体内部的电荷分布}
        \inputfigure[scale=0.75]{Chapter09A_Figure03}
    }\qquad
    \subfigure[导体表面的电荷分布]
    {
        \label{图:导体表面的电荷分布}
        \inputfigure[scale=0.75]{Chapter09A_Figure04}
    }
\end{Figure}
如\Refx{图}{导体内部的电荷分布}所示，在导体内任意一点处取一球面$\Sigma$作为高斯面，那么根据\Refx{定律}{电场高斯定理}，由于在导体内部的电场强度$\vb*{E}$处处为零
\begin{equation*}
    \frac{q}{\varepsilon_0}=\Isot{\vb*{E}\cdot\dif{\vb*{S}}}=0
\end{equation*}
因此球面$\Sigma$内的电荷量的代数和为零，由于球面可以取的任意小，这就可以说明，当静电平衡时，\textbf{导体内任意点处均无净电荷存在，导体所带电荷只能分布在导体的外表面上。}

从微观上看，当导体达到静电平衡时，净电荷的分布仅限于导体最外侧的一两个原子层内。

下面我们来探究导体表面的电荷分布与导体表面的电场强度间的关系。

如\Refx{图}{导体内部的电荷分布}所示，设导体表面某处的电荷面密度为$\sigma$，设该处的电场强度为$\vb*{E}$，以该处为中心，以该处的表面法线为轴线，构造一个小而扁平的圆柱形高斯面，圆柱面截导体表面的面积微元为$\dif{\vb*{S}}$，圆柱面的上下底面分别位于导体外部和内部。

对于位于导体内部的高斯面，由于导体内部的电场强度处处为零，因此其上的电通量为零。

对于位于导体外部的高斯面，由于导体表面的电场强度总是垂直于表面，因此，圆柱侧面上的电通量为零，圆柱上底面的电通量为$\dif\Phi_E=E\dif{\vb*{S}}$，而圆柱包围的电荷量为$\sigma\dif\vb*{S}$。

根据\Refx{定律}{电场高斯定理}
\begin{equation*}
    \frac{\sigma\dif{\vb*{S}}}{\varepsilon_0}=E\dif{\vb*{S}}
\end{equation*}
这就得到了
\begin{BoxEquation}[导体表面的电荷量与场强的关系]
    \sigma=E\varepsilon_0
\end{BoxEquation}
\Refx{公式}{导体表面的电荷量与场强的关系}指出，导体表面的场强与该处的电荷面密度成正比，即有$E=\frac{\sigma}{\varepsilon_0}$，但不能因此错误的认为，导体表面的场强是仅由该处的电荷$\sigma$产生的，因为根据\Refx{公式}{均匀无限大带电平面的电场}，由$\sigma$产生的电场应为$E=\frac{\sigma}{2\varepsilon_0}$，且分布于表面的两侧。但是在分布于导体表面的其他电荷的作用下，导体内侧的电场被抵消，导体外侧的电场被加强了一倍。这是导体的静电学特性带来的必然的电场分布结果。

\Refx{公式}{导体表面的电荷量与场强的关系}仅阐述了导体表面电荷分布与导体表面场强间的关系，并没有独立的说明表面电荷分布的规律，事实上，导体的表面电荷分布，不仅与导体自身的形状和带电量有关，同时与导体附近的外电场有关，是极为复杂的。

但是对于不受外电场作用的\uwave{孤立带电导体}，其电荷分布与表面性质具有以下大致关系
\begin{enumerate}
    \item 在导体表面较为凸出，即曲率为较大的正值的位置，电荷面密度较大。
    \item 在导体表面较为平坦，即曲率为零的位置，电荷面密度较小。
    \item 在导体表面较为凹陷，即曲率为较大的负值的位置，电荷面密度更小。
\end{enumerate}
即一般来说，\textbf{孤立带电导体某处的曲率越大，该处的电荷面密度就越大}。

为了验证这一说法，下面我们来看一个较为特殊的例子，设两个半径分别为$R_1$和$R_2$的导体球（$R_1>R_2$）相距很远，通过一根很长的细导线将两者连接。
\begin{itemize}
    \item 使两个导体球相距很远的目的是使其电场不会相互影响，从而可以看作孤立导体球。
    \item 使用细导线将两者的连接的目的是使其等电势，从而可以视为一个导体组。
\end{itemize}
现在令该导体组带上一定的电荷量，设两个导体球分别带电$q_1,q_2$，为求出$q_1,q_2$的关系，就要利用两者等电势的条件，但我们尚且还不清楚孤立带电导体球的电势和其电荷量的关系。

较为一般的，考虑一个带电量为$q$半径为$R$的孤立带电导体球，由于导体的性质，这些电荷量必然只能分布在导体球的表面，由于其是孤立的，为了使得导体内部电场强度处处为令，这些电荷量在其表面的分布必然是均匀的，这就可以视为一个均匀带电球壳的模型。

根据\Refx{公式}{带电球壳的电场分布}
\begin{equation*}
    E=\begin{cases}
        0~,&r<R\\[2mm]
        \dfrac{1}{4\pi\varepsilon_0}\dfrac{q}{r^2},&r>R
    \end{cases}
\end{equation*}
因此孤立带电导体球任何一处的电势，均可以用$R$至$+\infty$的过程中$\vb*{E}$的积分表示
\begin{equation*}
    V=\Int[R][+\infty]{\dfrac{1}{4\pi\varepsilon_0}\dfrac{q}{r^2}}\dif{r}=\dfrac{1}{4\pi\varepsilon_0}\dfrac{q}{R}
\end{equation*}
即
\begin{BoxEquation}[孤立导体球的电势]
    V=\dfrac{1}{4\pi\varepsilon_0}\dfrac{q}{R}
\end{BoxEquation}

回到该例子中，由于两个导体球等电势
\begin{equation*}
    \dfrac{1}{4\pi\varepsilon_0}\dfrac{q_1}{R_1}=\dfrac{1}{4\pi\varepsilon_0}\dfrac{q_2}{R_2}
\end{equation*}
从而
\begin{Equation}[孤立带电体电荷面密度1]
    \frac{q_1}{q_2}=\frac{R_1}{R_2}
\end{Equation}
式\ref{式:孤立带电体电荷面密度1}就说明两个等势的孤立导体球所带的电荷量，与导体球的半径成正比。

下面我们来考察两者的电荷面密度和半径的关系，由于孤立导体球的表面电荷均匀分布
\begin{equation*}
    q_1=4\pi R_1^2\sigma_1\qquad q_2=4\pi R_2^2\sigma_2
\end{equation*}
代入式\ref{式:孤立带电体电荷面密度1}得
\begin{Equation}[孤立带电体电荷面密度2]
    \frac{\sigma_1}{\sigma_2}=\frac{R_1^{-1}}{R_2^{-1}}
\end{Equation}
式\ref{式:孤立带电体电荷面密度2}就说明两个等势的孤立导体球所具有的电荷面密度，与导体球的半径成反比，即与导体球的曲率成正比，这就验证了“曲率越大，电荷面密度越大”的结论。

应指出的是，对于一般的导体，\textbf{曲率和电荷面密度间不存在明确的定量函数关系，曲率越大电荷面密度越大只是定性结论}，式\ref{式:孤立带电体电荷面密度2}的正比关系仅适用于此例的特殊情形。

以上结论的一个重要应用时，对于形状不规则的孤立带电导体，在其表面较为尖锐的地方，曲率极大，尖端表面的电荷面密度也会很大，尖端表面的电场强度会特别强，以至于有时足矣将空气分子电离，产生大量可以自由运动的阴阳离子，在尖端的强电场的作用下做激烈的加速运动，产生离子流，这种将空气击穿的放电现象，称为\uwave{尖端放电}。

尖端放电时，如果在尖端附近放置一个蜡烛，会明显观察到蜡烛的火焰（火焰本身就是由高温气体离子在电子跃迁中放出的光构成）在离子流的作用下倾斜，因此也形象的称为\uwave{电风}。

尖端放电在黑暗中观察时，会看到尖端附近笼罩了一层光晕，称为\uwave{电晕}，这是由尖端处的强电场的作用下产生的离子的电子在跃迁中放出的光辐射形成，与火焰的原理的类似，只不过前者的电离是由于电场作用，而后者的电力是由于氧化还原反应的作用。在夜间的高压输电线上常能观察到电晕的现象，这会导致电能的浪费，为此多采用表面光滑的粗导线。

尖端放电现象的一个重要应用是\uwave{避雷针}，当雷雨云中积聚了大量电荷时，在静电感应的作用下，避雷针作为地面上最为尖锐的导体，避雷针的尖端会具有最高的电荷面密度，从而空气最有可能在雷雨云和避雷针之间被击穿并放电，这就形成了闪电，避雷针可以将放电电流通过良好接地的导线安全引入大地。避雷针通过主动导引并消耗雷雨云中的电荷，从而避免电荷在雷雨云中进一步积累通过相对尖锐的高层建筑屋顶处击穿，造成雷击破坏建筑物。

\subsection{带有空腔的导体}
通过以上的讨论，我们已经对实心导体在静电场中的电荷分布和电场分布具有了一定了解，但是对于空心导体或者说带有\uwave{空腔}的导体，导体的内表面上是否有电荷？导体的空腔内是否有电场？如果空腔内还有电荷情况又会如何？这些问题的答案我们尚还不清楚。

为了讨论空腔导体，我们首先需要引入以下的一个准则。
\begin{BoxPrinciple}[静电场的唯一性原理]
    若空间中各个导体的形状大小和位置是已知的，此时，

    如果可以确定各个导体\textbf{所带的电荷量}或\textbf{所具有的电势}中的任意一个，

    那么，各个导体的电荷分布，与空间中的电场分布，就唯一的确定了。
\end{BoxPrinciple}
\Refx{定律}{静电场的唯一性原理}的严格证明需要涉及较为复杂的微分方程理论，这里从略。

\Refx{定律}{静电场的唯一性原理}的一个重要推论就是，对于某一确定的带电导体，如果我们能通过一些方式，首先得知该导体在某一电荷分布下，能使得其产生的电场符合要求，即导体中的电场处处为零，那么这种电荷分布就是其处于静电平衡时唯一可能的电荷分布。

下面，我们就将通过唯一性原理，研究在外电场中的带电空腔导体的电荷分布与电场分布。
\begin{Figure}{空腔导体的电荷分布与电场分布}
    \subfigure[实心导体]
    {
        \label{图:实心导体}
        \inputfigure[scale=0.75]{Chapter09A_Figure05}
    }\qquad
    \subfigure[空腔导体]
    {
        \label{图:空腔导体}
        \inputfigure[scale=0.75]{Chapter09A_Figure06}
    }
\end{Figure}
首先运用\Refx{定律}{电场高斯定理}，取一包围空腔但是完全位于导体内部的高斯面，由于导体内的电场强度处处为零，这一高斯面尚的电通量为零，这就可以推定导体内表面的电荷的代数和为零，但是电荷的代数和为零，存在以下两种可能的情况
\begin{enumerate}
    \item 导体的内表面上，处处不带电荷。
    \item 导体的内表面上，一部分带正电荷，一部分带负电荷，两者的代数和为零。
\end{enumerate}
对于一般的导体，高斯定理并不能确定究竟是哪一种情况，更进一步关于空腔内是否有电场的问题，高斯定理也不能回答，这就需要通过\Refx{定律}{静电场的唯一性原理}来分析。

如\Refx{图}{实心导体}所示，设想带电导体在未挖去空腔时，其电荷在外电场$\vb*{E}_0$的作用下达成了某种静电平衡，使得导体内部的电场处处为零。如\Refx{图}{实心导体}所示，在带电导体内部挖去空腔后，尽管我们并不知道现在的电荷分布，但是我们知道如果依照原先的电荷分布，此时仍然能使得相同的区域的电场处处为零，这一区域包含了现在的导体壳和导体内的空腔两部分，因此这一电荷分布可以使得空腔导体的导体部分电场处处为零。

那么根据\Refx{定律}{静电场的唯一性原理}，这就是空腔导体唯一可能的电荷分布。
\begin{BoxPrinciple}[空腔导体的电荷与电场分布]
    对于一个带电空腔导体，若空腔内没有电荷，那么电荷只能分布在导体的外表面上，

    同时，导体内表面处处没有电荷分布，导体空腔中电场强度处处为零。

    导体内是否有空腔，导体内空腔的大小和位置，均不会改变导体的电荷分布，

    导体外表面的电荷分布，仅由导体自身带电量和外电场的情况决定，与空腔无关。
\end{BoxPrinciple}

那么是否存在这样一种可能？当导体具有空腔时，其电荷实际上会采取一种与原先不同的分布，这种分布中外表面和内表面同时参与，这种分布仍然可以使得导体壳中的电场强度处处为零，但空腔内的电场强度则不再为零。\Refx{定律}{静电场的唯一性原理}指出这种情况是不存在的，因为电荷遵循原先的分布时已经可以使得导体壳内的场强处处为零，而给定条件下能使得导体内场强处处为零的电荷分布是唯一的，这就直接排除了其他情形。

\Refx{定律}{空腔导体的电荷与电场分布}指出，无论导体壳是否带有电荷，无论导体外的电场多么复杂，空腔内的电场强度恒为零，这在电力工程中是有应用的，在从事高压带电作业时工作人员会全身穿戴金属丝编织的防护衣帽、手套、鞋子，这就相当于一个将工作人员包裹在空腔内的导体壳，这就能屏蔽外部高压电的强电场从而保护工作人员的安全。

\Refx{定律}{空腔导体的电荷与电场分布}讨论了空腔内没有电荷的情况，我们知道此时外电场对空腔是无影响的，但是如果空腔内有电荷，此时我们会思考两个问题
\begin{itemize}
    \item 导体能否对内屏蔽外部电场（正向屏蔽），即空腔内的电场是否会受到外电场的影响。
    \item 导体能否对外屏蔽内部电场（逆向屏蔽），即空腔内的电荷分布是否会改变外电场。
\end{itemize}
这两个问题在生产实践中均是有应用的
\begin{itemize}
    \item 通过正向屏蔽，我们可以保护精密的电器设备不受外部电场的干扰和破坏。
    \item 通过逆向屏蔽，我们可以对外界屏蔽高压器件的电场从而避免发生意外。
\end{itemize}
这两个问题的回答取决于导体是否接地，我们分别予以讨论。

1. 如\Refx{图}{导体空腔接地的情况}所示，讨论空腔内有电荷且导体接地的情况。
\begin{Figure}{导体空腔接地的情况}
    \inputfigure[scale=0.52]{Chapter09A_Figure07}
\end{Figure}
在子图[3]中，导体壳接地，导体壳外部具有电场$\vb*{E}_0$，导体的空腔内存在电荷$q_2$。

在子图[3]中，根据\Refx{定律}{电场高斯定理}，如蓝色虚线所示，取一包围空腔却完全位于导体内的高斯面，由于导体内电场强度$\vb*{E}_b=\vb*{0}$，因此高斯面内电荷的代数和为零，导体中不可能有电荷，而在导体空腔内已知具有$q_2$的电荷，故在导体内表面上应有$-q_2$的电荷。

这就得出了一个重要结论\footnote{通过该结论的推导过程可以看出，该结论对于接地或不接地的导体都是成立的。}
\begin{BoxPrinciple}[导体空腔和内表面的电荷量关系]
    导体内表面的电荷量，总是与导体空腔内的电荷量等量异号。
\end{BoxPrinciple}
那么现在我们要研究的两个问题就是
\begin{enumerate}
    \item 接地导体能否屏蔽外部电场对内的影响，为了验证这一点，我们移除子图[3]中的外部电场$\vb*{E}_0$得到子图[2]，并判断两者在导体空腔的电场强度$\vb*{E}_c$是否完全相同。
    \item 接地导体能否屏蔽内部电场对外的影响，为了验证这一点，我们移除子图[3]中的空腔电荷$q_2$\hspace{6pt}得到子图[1]，并判断两者在导体外部的电场强度$\vb*{E}_a$是否完全相同。
\end{enumerate}
在子图[1]中\footnote{应注意，导体接地与导体是否带电完全无关，导体接地仅代表导体电势为零，故这里导体壳带电量未必为$0$。}，设接地导体壳在外电场$\vb*{E}_0$下具有电荷$q_1$，根据\Refx{定律}{空腔导体的电荷与电场分布}，此时导体空腔内的电场强度$\vb*{E}_c=\vb*{0}$，且导体内表面没有电荷分布，故导体外表面的电荷量为$q_1$，设其呈$(q_1)$分布，以$\vb*{G}$表示其场强\footnote{G-Ground，代表接地，这里是指接地的导体壳产生的场强。}，故导体外的电场强度$\vb*{E}_a=\vb*{E}_0+\vb*{G}$。

在子图[2]中，导体空腔内的电荷量为$q_2$，设其以$(q_2)$分布，根据\Refx{定律}{导体空腔和内表面的电荷量关系}，导体内表面的电荷量为$-q_2$，设其以$(-q_2)$分布。

由于导体壳接地电势为零，且此时没有外电场的存在，因此该导体壳是空间中唯一的导体，而以下引理指出\footnote{引理的证明从略。}，\textbf{若所有导体的电势为零，则导体以外空间的电势处处为零}，电场强度是电势的负梯度，电势处处为零，电场强度也处处为零，故导体外的电场强度$\vb*{E}_a=\vb*{0}$。

由于导体外的电场强度处处为零，如子图[3]的红色虚线所示，取一包围导体的高斯面，根据\Refx{定律}{电场高斯定理}，该高斯面内的电荷量的代数和为零，而导体空腔内的$q_2$与导体内表面上的$-q_2$已是正负抵消的，因此导体外表面的总电荷量为$0$，但这不代表导体外表面处处没有电荷，而是有正有负相互抵消的，记外表面的电荷分布为$(0)^{*}$。

以$\vb*{C}$表示$(q_2),(-q_2),(0)^{*}$在空腔中产生的场强之和\footnote{C-Cavity，代表空腔，这里是指空腔内产生的场强。}，故导体空腔内的电场强度$\vb*{E}_c=\vb*{C}$。

而子图[3]可以视为子图[1]和子图[2]的叠加，由于电荷和电场在叠加后，仍然能使得导体内部$\vb*{E}_b=\vb*{0}$，根据\Refx{定律}{静电场的唯一性原理}可知，这一电荷电场分布是“导体壳接地，导体空腔内有电荷$q_2$，导体外有外电场$\vb*{E}$”时唯一可能的情况，因此，导体外表面的电荷分布为$(q_1)+(0)^{*}$，导体壳带电$q_1-q_2$，导体壳外$\vb*{E}_a=\vb*{E}_0+\vb*{G}$，导体空腔内$\vb*{E}_c=\vb*{C}$。
\begin{enumerate}
    \item 对比子图[3]和子图[2]，两者在导体空腔的电场强度$\vb*{E}_c$相等，均为$\vb*{E}_c=\vb*{C}$。
    \item 对比子图[3]和子图[1]，两者在导体外部的电场强度$\vb*{E}_a$相等，均为$\vb*{E}_c=\vb*{E}_0+\vb*{G}$。
\end{enumerate}
综上，我们可以总结出以下的结论。
\begin{BoxPrinciple}[空腔接地导体的基本性质]
    对于一个接地的空腔导体，其具有双向的电场屏蔽效应，

    腔内电场不受腔外电场和电荷的影响，

    腔外电场不受腔内电场和电荷的影响，

    导体外表面的电荷总量由外电场的情况决定，腔内电荷仅能影响外表面电荷分布，而不改变总量。

    导体内表面的电荷总量与腔内电荷等量异号，分布仅由腔内电荷的位置决定，与外电场无关。
\end{BoxPrinciple}
\Refx{定律}{空腔接地导体的基本性质}所描绘的现象称为\uwave{静电屏蔽}。

2. 如\Refx{图}{导体空腔未接地的情况}所示，讨论空腔内有电荷且导体未接地的情况。
\begin{Figure}{导体空腔未接地的情况}
    \inputfigure[scale=0.52]{Chapter09A_Figure08}\vspace{-15pt}
\end{Figure}
在子图[1.b]和子图[3.b]中，导体壳均带有电荷量$q_1$，导体空腔内包含电荷$q_2$，导体壳均未接地，两者的差异在于，前者外部无电场，后者外部具有$\vb*{E}_0$的外电场。

在子图[1.b]和子图[3.b]中，根据\Refx{定律}{导体空腔和内表面的电荷量关系}和\Refx{定律}{电荷守恒定律}，两者的内表面的电荷量均为$-q_2$，两者外表面的电荷量均为$q_1-q_2$。

那么现在的问题就是，子图[1.b]和子图[3.b]在导体空腔内的电场强度$\vb*{E}_c$是否相同？如果其相同，就说明导体壳即便没有接地，也可以对空腔内部屏蔽外部的电场。

为了回答这一问题，我们构造子图[1.a],[2],[3.a]
，使得[1.a]$+$[2]$=$[1.b]和[3.a]$+$[2]$=$[3.b]。

\textbf{对于子图[1.a]和子图[3.a]}，从结果上看，其相当于在子图[1.b]和子图[3.b]的基础上，移除导体空腔的电荷$q_2$和伴随之的导体内表面的电荷$-q_2$，但仍使外表面带电$q_1+q_2$。

从构造上看，导体壳均带有电荷量$q_1+q_2$，导体空腔内均无电荷，根据\Refx{定律}{空腔导体的电荷与电场分布}，导体内表面处处无电荷，导体空腔内电场强度处处为令，即$\vb*{E}_c=0$，故外表面带电量$q_1+q_2$，设无外电场时以$(q_1+q_2)$分布，设有外电场时以$(q_1+q_2)'$分布。

设导体壳在无外电场和有外电场时，导体外表面带单位电荷时产生的场强\footnote{S-Surface，代表表面，这里指外表面带单位电荷时产权的场强。}分别为$\vb*{S}$和$\vb*{S}^{*}$，因此在子图[1.a]中场强$\vb*{E}_a=\vb*{S}(q_1+q_2)$，而在子图[3.a]中场强$\vb*{E}_a=\vb*{S}^{*}(q_1+q_2)+\vb*{E}_0$。

\textbf{对于子图[2]}，从结果上看，其相当于在子图[1.b]的基础上使得导体外表面不带电。

从构造上看，导体壳带有的电荷量为$-q_2$，导体空腔内电荷量为$q_2$，根据\Refx{定律}{空腔导体的电荷与电场分布}，导体内表面带电$-q_2$，导体外表面不带电，设其以$(0)^{*}$分布，设由空腔电荷和导体内外表面$(q_2),(-q_2),(0)^{*}$产生的合场强为$\vb*{C}$，故$\vb*{E}_c=\vb*{C}$，而对于导体壳外的电场，尽管子图[2]中的导体壳没有接地，但是其与\Refx{图}{导体空腔接地的情况}中接地的子图[2]具有完全相同的电荷分布，因此必然产生完全相同的电场，故$\vb*{E}_a=\vb*{0}$。

通过[1.a],[2]叠加和[3.a],[2]叠加，应用\Refx{定律}{静电场的唯一性原理}可知，对于[1.b]和[3.b]，导体外表面电场分布分别为$(q_1+q_2)+(0)^{*}$和$(q_1+q_2)'+(0)^{*}$，导体壳外的电场分别为$\vb*{E}_a=\vb*{S}(q_1+q_2)$和$\vb*{E}_a=\vb*{S}'(q_1+q_2)+\vb*{E}_0$，导体空腔内电场均为$\vb*{E}_c=\vb*{C}$。

这就说明，\textbf{对于未接地的空腔导体，空腔内部的电场不受外部的影响}。

那么反之是否是成立呢？答案是否定的，欲研究空腔外部的电场是否受内部的干扰，就应在保持子图[1.b]和[3.b]外壳带电量为$q_1$时，移除空腔内的电荷$q_2$并考察$\vb*{E}_a$是否变化。

这会得到子图[1.c]和子图[3.c]，由于这里又回到空腔中无电荷的情形，因此不难看出，两者导体壳外电场分别为$\vb*{E}_a=\vb*{S}q_1$和$\vb*{E}_a=\vb*{S}'q_1+\vb*{E}_0$，但这显然相较于原先发生了变化。

这就说明，\textbf{对于未接地的空腔导体，空腔外部的电场会受内部的影响}。

但是，这种内对外的影响是否是完全随意的？从以下两个角度考虑
\begin{enumerate}
    \item 我们将子图[3.b]中的$q_2$移动一些位置，这会使我们得到三个与原先不同的电荷分布$(q_2),(-q_2),(0)^{*}$，因此，导体外表面的电荷分布$(q_1+q_2)'+(0)'$将会随之发生变化，但是，导体外的场强$\vb*{E}_a=\vb*{S}'(q_1+q_2)+\vb*{E}_0$并不会发生改变，这就表明，\textbf{对于未接地的空腔导体，在空腔内电荷量一定时，改变电荷的位置不会改变导体外的电场。}
    \item 我们关注到子图[3.a]可以视为子图[3.b]中将空腔内的电荷$q_2$完全转移到导体壳上的结果，两者的导体外表面电荷量均为$q_1+q_2$，但具有不同的电荷分布，然而导体外的场强$\vb*{E}_a=\vb*{S}'(q_1+q_2)+\vb*{E}_0$却是完全相同的，这就表明，\textbf{对于未接地的空腔导体，其产生的场强仅与导体和空腔内部的电荷总量有关，而与这些电荷的具体分配无关。}
\end{enumerate}
综上，我们可以总结出以下的结论。
\begin{BoxPrinciple}[空腔未接地导体的基本性质]
    对于一个未接地的空腔导体，其仅有单向的电场屏蔽效应，
    
    腔内电场不受腔外电场和电荷的影响，
    
    腔外电场会受腔内电场和电荷的影响，
    
    但是这种影响是有限的，
    
    空腔内的电荷量一定时，改变电荷位置不会改变导体外的电场。

    空腔导体在外部产生的电场强度，仅与导体和空腔内部的电荷总量有关，与电荷的具体分配无关。

    也就是说，将空腔内部的电荷完全转移到导体壳上，不会改变外部的电场。
\end{BoxPrinciple}



